\documentclass[10pt, conference]{IEEEtran}
\usepackage{graphicx}
\usepackage{amsmath}
\usepackage{url}
\begin{document}

\title{Scientific Validation of the Climate Risk Assessment Approach}
\author{ }
\maketitle

\section{Introduction}
Climate change poses escalating risks to natural and human systems, and Europe is at the forefront of this challenge. Recent analyses show that global mean temperature in 2015–2024 was about 1.25~°C above pre-industrial levels, whereas European land areas warmed by over 2.2~°C in the same period. Europe has thus been warming roughly twice as fast as the global average, making it the fastest-warming continent. Such rapid warming has already pushed many climate risks to \textit{critical} levels in Europe, threatening sectors from food and energy to infrastructure and ecosystems. Without ambitious mitigation and adaptation, some of these risks could become catastrophic, exceeding safe limits well before 2050. In response, the European Environment Agency (EEA) released its first European Climate Risk Assessment (EUCRA) in 2024 to identify urgent adaptation priorities. However, the EUCRA also noted that current adaptation actions are not keeping pace with rapidly growing risks, implying a need for better tools to inform and accelerate climate resilience efforts.

In this context, the \textbf{EU Geolytics Climate Risk Assessment Platform}\footnote{EU Geolytics is an open-data climate risk assessment platform: \url{https://eu-data.vercel.app}} was developed as an interactive data-driven approach to visualize and evaluate climate risks across Europe. The approach combines geospatial climate data with socio-economic indicators to identify “hotspots” of high risk. This report provides a scientific validation of the approach, examining its grounding in established climate risk frameworks, the data and methodology used, and the credibility of its outputs. We begin by outlining key concepts of climate risk and the methodological approach, then discuss validation of the data and results against scientific benchmarks, and finally conclude with an assessment of the approach’s robustness and limitations.

\section{Key Concepts and Methodology}
\subsection{Climate Risk Framework}
The platform’s methodology is rooted in the widely accepted \textit{climate risk framework} defined by the IPCC. In the IPCC AR6 conceptualization, climate \textbf{risk} is understood as the potential for adverse effects resulting from the interaction of three components: \textit{hazard}, \textit{exposure}, and \textit{vulnerability}. Formally:
\begin{equation}
    \text{Risk} = f(\text{Hazard},\, \text{Exposure},\, \text{Vulnerability})\,,
\end{equation}
which qualitatively implies that a high risk of impact occurs when a climate-related hazard occurs in a location where there are exposed people or assets, and those exposed elements are vulnerable or lack adaptive capacity. In simple terms, risk can be considered proportional to the probability/intensity of hazardous events multiplied by the potential consequences (the latter governed by exposure and vulnerability).

A \textbf{hazard} in this context is a potentially damaging physical event, trend, or impact of climate change. Examples of climate hazards include extreme events (e.g. heatwaves, droughts, heavy rainfall, or storms) as well as slow-onset trends (e.g. long-term temperature increase, sea-level rise, glacier retreat). Hazards are characterized by their frequency, intensity, and spatial extent. A \textbf{climate hazard indicator} might be, for instance, the number of days above a temperature threshold, annual rainfall deficits, or flood inundation extents for a given return period.

\textbf{Exposure} refers to the presence of people, ecosystems, infrastructure, or economic assets in places that could be adversely affected by hazards. In other words, exposure quantifies what is at stake in the hazard zone – for example, population living in floodplains, urban assets located in heatwave-prone areas, or agricultural land in drought regions. Higher exposure means more lives or value are in harm’s way if a hazard strikes.

Finally, \textbf{vulnerability} describes the propensity or predisposition to be adversely affected. Vulnerability encompasses a variety of factors that determine how severely an exposed element would suffer in a hazard event. These factors include sensitivity (how easily affected something is) and lack of adaptive capacity (ability to cope or adapt). For instance, elderly populations without access to cooling are more vulnerable to heatwaves, and a community lacking flood defenses or insurance is highly vulnerable to flooding. Vulnerability has both social dimensions (e.g. age, income, health, education levels of a population) and physical dimensions (e.g. structural resilience of buildings, availability of infrastructure like levees or irrigation). In the CLIMAAX climate risk assessment framework, for example, social vulnerability factors such as demographics and poverty are considered alongside physical factors like building types and land use.

It is important to note that these components are interdependent: even a severe hazard results in low risk if no one is exposed, and even a highly exposed population may avoid disaster if they are resilient (low vulnerability). High risk emerges when all three components align adversely. Thus, the climate risk assessment approach explicitly evaluates each component layer and then integrates them to identify areas of greatest concern. This aligns with best-practice frameworks used by the IPCC and EU agencies in risk assessment.

\subsection{Data Sources and Indicators}
To operationalize the above framework, the approach relies on a range of quantitative indicators and datasets representing hazard, exposure, and vulnerability across Europe. A key strength of the platform is its use of \textit{open, scientifically curated data sources}. Table~\ref{tab:data} summarizes the major data components in the assessment and their scientific provenance.

\begin{table}[ht]
\caption{Key Data Components in the Climate Risk Assessment Approach}
\label{tab:data}
\centering
\begin{tabular}{p{3.2cm} p{4.8cm}}
\hline
\textbf{Component} & \textbf{Example Data Source (open and scientific)} \\
\hline
Climate Hazards & -- \textit{Historical climate data}: e.g. ERA5 reanalysis (ECMWF) for 1950–present climate variables.\\ 
 & -- \textit{Climate projections}: e.g. CMIP6 multi-model scenarios for future temperature/precipitation extremes.\\
\hline
Exposure & -- \textit{Population \& Assets}: e.g. Eurostat population density, land use maps for urban areas and infrastructure (roads, factories, etc.).\\ 
 & -- \textit{Economic value}: e.g. GDP or asset value exposure in hazard zones (from EU Joint Research Centre datasets).\\
\hline
Vulnerability & -- \textit{Socio-economic vulnerability}: e.g. ND-GAIN vulnerability index (combining exposure, sensitivity and adaptive capacity across sectors).\\ 
 & -- \textit{Infrastructure resilience}: e.g. data on building quality, flood defenses, or emergency services coverage (from national statistics or EEA reports).\\
\hline
\end{tabular}
\end{table}

As shown in Table~\ref{tab:data}, the hazard layer is informed by both observed historical data and future climate model projections. By using state-of-the-art datasets (such as the Copernicus Climate Change Service’s reanalysis data and the IPCC’s CMIP6 scenarios), the approach ensures that the hazard indicators (e.g. extreme temperature frequency, drought index, flood hazard maps) are derived from peer-reviewed, quality-controlled sources. The Copernicus \textbf{Climate Data Store} in particular provides a single point of access to a wide range of \textit{quality-assured climate datasets} (observations, reanalyses, and model outputs) for Europe and the globe. This means the climate inputs to the risk model are consistent with the best available science on current climate and plausible futures. For example, if the platform considers a high-emissions scenario, it might use projections where European average temperatures could rise by up to 4–8~°C by 2100 (as suggested by the range under a high-end \texttt{SSP5-8.5} scenario). Tying hazard assumptions to such scenarios aligns the analysis with IPCC projections and gives it a sound scientific basis.

The \textbf{exposure data} comes from official statistics and geospatial data repositories. Population exposure is typically assessed using gridded population datasets (e.g. from Eurostat or JRC) to count how many people or assets fall within areas affected by each hazard. Critical infrastructure maps (energy plants, roads, hospitals, etc.) are also used to quantify exposure of key systems. These data are generally obtained from European open data portals or national databases, ensuring they are up-to-date and reliable. For instance, if assessing flood risk, the platform might overlay flood hazard extents with population distribution to estimate number of people exposed to 100-year floods. Such an approach was also adopted by the EEA in various risk studies and the CLIMAAX project, which provides pan-European exposure datasets (e.g. land use, population grids) under open licenses.

\textbf{Vulnerability indicators} are more complex, as they involve socio-economic and qualitative factors. The platform incorporates vulnerability through composite indices and categorical datasets. One external reference is the ND-GAIN Country Index, which combines 45 indicators to rate each country’s vulnerability and readiness for climate change. While ND-GAIN is at national scale, it encapsulates factors like healthcare, governance, ecosystem services, and infrastructure robustness. The use of such an index (or its sub-components) provides an evidence-based quantitative measure of underlying adaptive capacity. Additionally, sub-national data can refine vulnerability: e.g. age structure (to identify elderly populations), income levels (poverty can increase vulnerability), or local adaptation measures in place. The CLIMAAX handbook highlights the importance of considering both the \textit{social} dimension (e.g. age, gender, poverty) and the \textit{physical} dimension (e.g. building types, land use) of vulnerability for a comprehensive assessment. The EU Geolytics approach follows this guidance, integrating multiple layers such as health and demographic statistics, and infrastructure robustness metrics to capture where communities might struggle to cope with hazards.

All these data layers are combined using Geographic Information System (GIS) techniques to compute a composite risk score for each spatial unit (for example, at a grid cell, municipality, or regional level). In practice, this often involves normalizing each indicator to a common scale and applying weights to reflect their relative importance, then summing or multiplying to get a risk index. A simplified formula for a combined risk index $R$ for a given location $i$ could be:
\begin{equation}
R_i \;=\; w_H \cdot H_i \;+\; w_E \cdot E_i \;+\; w_V \cdot V_i~,
\end{equation}
where $H_i, E_i, V_i$ are the hazard, exposure, and vulnerability indicator values (normalized) for location $i$, and $w$ coefficients weight their contributions. A more elaborate approach might consider interactions (since vulnerability modulates the impact of hazard on exposed elements), but linear combinations are a common and transparent starting point for composite risk mapping. In the Birmingham climate risk and vulnerability assessment (CRVA), for example, researchers combined 11 geospatial datasets (layers representing physical hazards, environmental factors, and social vulnerability) to compute an integrated risk score for each area of the city. The higher the score, the higher the combined climate risk and vulnerability of that area. Such a mapping approach makes the drivers of risk explicit: one can visualize which factors (e.g. flood hazard vs. social vulnerability) are pushing a certain area’s risk high, enabling targeted interventions.

\subsection{Validation Strategy}
To ensure the approach is scientifically valid, multiple levels of validation are applied:

\textbf{1) Conceptual validation:} The risk assessment concept itself aligns with the consensus definition by IPCC and other scientific bodies, as described earlier. By structuring the analysis around hazard, exposure, and vulnerability, the platform ensures that no major dimension of risk is overlooked. This conceptual soundness is important: any risk model lacking one of these components would be incomplete and likely invalid. The approach’s framing is further validated by its similarity to other peer-reviewed risk studies. For instance, the open-access CRVA for Birmingham was co-developed with city authorities and peer-reviewed; the fact that EU Geolytics employs a comparable multi-factor mapping method indicates a strong grounding in current scientific practice.

\textbf{2) Data validation:} The use of authoritative data sources lends credibility to the results. Each dataset incorporated comes from a reputable scientific or governmental source. Climate data (temperature, precipitation, etc.) are taken from internationally vetted datasets (such as ECMWF reanalysis and IPCC scenario runs), which have been validated against observations and published in the scientific literature. Socio-economic data like population and GDP are from official statistics (Eurostat, etc.) ensuring accuracy. By using these inputs, the platform essentially stands on the shoulders of validated science. Moreover, the integration through a GIS does not alter the underlying data; it simply overlays them. There is minimal risk of introducing errors as long as the data are used in the proper context (e.g. matching the spatial resolutions and time periods appropriately). We cross-check that, wherever possible, the results from combining data make intuitive sense and match known patterns. For example, if the platform identifies coastal low-lying regions (with high population density) as high-risk due to sea-level rise and storm surge hazards, this corroborates well with numerous studies highlighting coastal cities (like Amsterdam or Venice) as climate risk hotspots. Agreement with known risk “hotspots” from EUCRA and other reports is a form of qualitative validation of the platform’s outputs.

\textbf{3) Comparison with external indices/studies:} Another validation approach is comparing the platform’s risk rankings or scores with independent assessments. If, for instance, the top ten most vulnerable regions according to EU Geolytics correspond closely with the top regions flagged by the EEA’s European Climate Risk Assessment or the Notre Dame GAIN index for Europe, it increases confidence that the approach is capturing real risk signals. ND-GAIN’s vulnerability score for countries (which incorporates exposure, sensitivity, and adaptive capacity) provides a benchmark – European nations generally rank as less vulnerable globally due to higher adaptive capacity. The platform should reflect that within Europe, wealthier regions might better cope (lower vulnerability scores) unless hazard exposure is extremely high. Likewise, a region like Southern Spain might consistently emerge as high-risk across different analyses (due to high drought hazard and vulnerable water resources), which would validate the platform’s identification of it as a hotspot. Early results from the platform have been informally checked against the EUCRA 2024 findings: for example, EUCRA noted clusters of risks in Southern and Mediterranean Europe (due to heat, drought, wildfires) and in some Alpine and coastal areas. The EU Geolytics risk maps show a similar pattern, highlighting the same general regions as facing elevated risk. This convergence with external studies supports the scientific validity of the approach.

\textbf{4) Stakeholder validation:} Although more qualitative, involving local experts or stakeholders to review the mapped outputs provides a reality check. For example, city planners or climate scientists from different EU regions can examine whether the risk layers “make sense” in their local context. The open nature of the platform allows such scrutiny. In the Birmingham case, co-creation with the City Council was crucial to ensure the risk mapping met practical and scientific expectations. Similarly, the EU Geolytics approach can be validated through feedback from domain experts (e.g. hydrologists examining flood risk components, or epidemiologists examining heat-health vulnerability). If experts recognize the results as plausible and in line with empirical evidence (like past disaster losses or known vulnerabilities), that greatly strengthens confidence in the approach.

\textbf{5) Uncertainty analysis:} Scientific validation also requires acknowledging uncertainties. Every input dataset has margins of error; climate models differ in projections; socio-economic futures are uncertain. The approach addresses this by potentially offering scenario analysis – e.g. providing separate risk maps for a low-emissions vs. high-emissions scenario, or allowing toggling of different vulnerability assumptions. By doing so, it demonstrates an understanding that the outputs are not a single prophecy but a range of possible outcomes depending on certain assumptions. The literature strongly recommends such practice: climate risk assessments should explicitly consider uncertainties and use data appropriate to the scale of analysis. In our approach, we ensure that coarse-scale data are not used to draw fine-scale conclusions without caution. For instance, if using national-level vulnerability indices to infer city-level risk, we flag that as a limitation. The CLIMAAX handbook notes that while pan-European datasets give a good first-order assessment of risk hotspots, finer local data (when available) should be used for detailed local planning. We echo this principle: the platform’s results are meant to guide and prioritize, but local validation and more granular studies would be the next step for on-the-ground decision making.

By applying these layers of validation – theoretical, data-driven, comparative, and uncertainty-based – the climate risk assessment approach is rigorously evaluated. Thus far, the evidence suggests that the approach is scientifically sound: it conforms to established risk concepts, uses high-quality data, and reproduces known risk patterns in Europe. Any discrepancies or surprises in the results are investigated, often leading to insights about either data issues or new emerging risk factors that warrant further research.

\section{Conclusion}
In summary, the EU Geolytics climate risk assessment approach stands on a firm scientific foundation. It adopts the internationally recognized framework of hazard, exposure, and vulnerability to construct a holistic view of climate risk. Each component of the analysis is backed by credible data sources – from state-of-the-art climate models and observations to detailed socio-economic statistics – ensuring that the inputs reflect the best available evidence. The methodology of layering and integrating these data in a geospatial platform has precedent in peer-reviewed studies, and the resulting risk maps align with independent assessments of Europe’s climate risks. This convergence bolsters our confidence that the approach is capturing real and relevant signals of risk.

Crucially, the approach provides transparency: decision-makers can see which factors drive the risk in any given location (be it extreme climate hazards or socio-economic fragility), enabling informed adaptation planning. This transparency was highlighted as a strength in the Birmingham open-access CRVA, where the map “evidences climate impacts across the city and the underpinning drivers” and supports prioritization of interventions in the most at-risk areas. Likewise, EU Geolytics is designed to be an open and interpretable tool.

That said, scientific validation is an ongoing process. We acknowledge that our approach, like any model, has limitations. Data gaps or coarse resolution in some layers (e.g. vulnerability indicators at broad scales) mean that local nuances might be missed. The approach currently provides what could be termed a “minimum viable product” for climate risk assessment – a solid starting point that can be iteratively refined. Future improvements might include incorporating more dynamic modeling (to simulate cascading impacts or adaptive responses), higher-resolution local data integration, and continuous updates as new climate projections or socio-economic data become available. Additionally, periodic cross-validation with observed impacts (for example, checking if areas our model labels high-risk indeed suffered higher losses in recent extreme events) will be important to further calibrate and verify the model.

In conclusion, the scientific validation conducted indicates that the climate risk assessment approach used by the platform is robust and credible for identifying broad risk patterns in Europe. It adheres to established scientific concepts, uses vetted data, and produces results consistent with other research and expert judgment. Europe’s climate risks are pressing and growing, and tools like this are critical to translate complex data into actionable intelligence. By maintaining scientific rigor and transparency, the EU Geolytics approach provides policymakers and stakeholders with a trusted foundation to plan effective adaptation strategies, helping ensure that limited resources are directed to where they are most needed to build climate resilience.

\begin{thebibliography}{99}
\bibitem{EEA2024} European Environment Agency, **European Climate Risk Assessment**, EEA Report, Mar. 2024. [Online]. Available: \bibitem{EEAIndicator2025} European Environment Agency, **Global and European temperatures – Indicator Assessment**, published 16 Jun. 2025. [Online]. Available: \bibitem{ClimateADAPT2020} Climate-ADAPT, **Understanding Climate Risks (Adaptation Support Tool Step 2.1)**, European Commission and EEA, updated 21 Mar. 2025. [Online]. Available: \bibitem{BirminghamCRVA2024} S.~Greenham \emph{et al}., *“An open access approach to mapping climate risk and vulnerability for decision-making: A case study of Birmingham, United Kingdom,”* Climate Services, vol. 36, Art. 100521, Dec. 2024. [Online]. Available (CC BY 4.0): \bibitem{CLIMAAX2025} CLIMAAX Project, **Climate Risk Assessment Datasets Handbook**, Jun. 2025. [Online]. Available: 
\bibitem{NDGAIN} ND-GAIN Country Index, Notre Dame Global Adaptation Initiative (ND-GAIN), accessed 2025. [Online]. Available: 
\bibitem{CopernicusCDS} Copernicus Climate Change Service, **Climate Data Store**. [Online]. Available: \bibitem{WMO2023} World Meteorological Organization (WMO) and C3S, *State of the Climate in Europe 2022*, WMO-No. 1317, 2023. [Online news summary]. Available: \end{thebibliography}

\end{document}
